\begin{exercise}[The zero-momentum vacuum-polarization identity]{I-CH09-013}
In the Schwinger representation of the one-loop photon two-point function,
keep the factor $s^{1-d_{\mathrm{reg}}/2}$ supplied by
$\dd\alpha\,\dd\beta=s\,\dd s\,\dd x$ and the Gaussian loop integral.  At
$p^2=0$ the right-hand side of the equation for
$(d_{\mathrm{reg}}-1)p^2\Pi(p^2)$ is proportional to
\[
  \int_0^\infty\dd s\,s^{1-d_{\mathrm{reg}}/2}\ee^{-sm_0^2}
  \left[-d_{\mathrm{reg}}m_0^2
  +\frac{d_{\mathrm{reg}}(2-d_{\mathrm{reg}})}{2s}\right].
\]
Show in a convergence strip that this is an integral of a total derivative,
then continue the identity dimensionally.  Deduce that the polarization tensor
has no momentum-independent photon-mass term and that $\Pi(p^2)$ has no
$1/p^2$ pole for $m_0>0$.  State how the massless limit is defined and check
the boundary terms.
\end{exercise}
